\section{Discussion}
Pairing the candidate and anchor critical values makes the uncertainty of a borrowing gain depend on how much is borrowed. The covariance-set result gives a finite-sample benchmark, the feasible multiplier construction retains the cancellation at the local reference-information scale, and the weak-signal experiment relates the reference information needed to certify a correction to the squared size of its attainable gain.

The matched ablation identifies the contribution of pairing. A small simultaneous gain can be estimable as a difference while separate bounds for its two terms remain too wide, and the dense-grid and non-Gaussian experiments show that this mechanism operates beyond a short threshold vector. Numerical integration and statistical calibration remain separate components of uncertainty. The oracle comparison is made on its stated certification event, and the implemented statistical guarantee is asymptotic; a limitation of the present implementation is that finite-sample coverage in scarce neighbourhoods must be read from the reported simulations rather than from the theorem.

The candidate class sets the remaining limits. Matrix projection can exploit a reversed relationship outside the nonnegative scalar path, and Full can be more efficient under strong alignment; Paired controls whether the fitted correction is used along the proposed path. The complete comparisons report these costs alongside the gains.

In the air-monitoring study, historical mean calibration and subsequent distributional verification answer different questions. A large lower-concentration discrepancy is established with a modest audit, whereas a one-percentage-point upper-tail difference may require nearly complete verification under the same simultaneous family. Designed reference sampling separates this precision requirement from environmental population sampling.

\section*{Data availability}
The EPA analysis table, source-verified date mapping and original computational archive are available in the public repository \url{https://github.com/Stork343/surrogate-assisted-local-conditional-distribution-inference}, source commit \texttt{70c47c512ab485edf502f16914f9ff481cc6be88}. The V3 reproduction package dated 29 September 2026 accompanying this manuscript contains the revised code, all generated results, input checksums and commands for rebuilding its tables and figures. The processed audit population and the historical calibration sample are identified separately.

\section*{Acknowledgements and funding}
This work was partially supported by the Beijing Natural Science Foundation (1242005), the Fundamental Research Funds for the Central Universities and the Research Funds of Renmin University of China (25XNN015), and the Ministry of Education Humanities and Social Sciences Research General Project (25YJA910005).

ChatGPT assisted the methodological revision, mathematical derivations, software development, numerical checks and drafting. Earlier computational development also used OpenAI Codex. These tools are not authors; responsibility for the submitted research, analysis and text rests with the authors.

\section*{Conflict of interest}
The authors declare no conflicts of interest.

\section*{Supplementary material}
The Supplement provides proofs, primitive and growing-grid conditions, numerical calibration details, rotating-role inference, the design-based audit extension and complete numerical results. The reproduction package provides executable code and the input and output data.

